The proposed pipeline will progress from auditable observations and frozen representations to parallel input- and loss-evidence channels, matched experimental comparisons, and uncertainty-qualified research outputs. The independent synthetic experiment will be the primary evidence source because it supplies known generating mechanisms and bounded risk truth. Dependent controls will test robustness; natural case studies will examine descriptive transfer without inventing population-harm labels. The twelve steps below specify proposed work and the assumptions under which its conclusions would be interpretable \cite{cerqueira,window,souza}.

\subsection*{Step 1 --- Define the Research Input}
The input will be an ordered tabular stream, a prediction log and a separately scheduled label-arrival log. A logistic classifier and a shallow decision-tree classifier will be fitted before evaluation and frozen. Index $i$ will refer only to test-stream observations: input time will be $t_i=i$ and label arrival will be $r_i=i+db_0$, where $b_0=1{,}000$ is the primary cohort size and $d$ the delay in primary-cohort units. Each prediction will precede release of its label. At test time zero, the historical model-training and risk-reference labels will already be fully resolved; the release schedule will apply only to test labels. This fully labelled warm start will make the risk reference available for the first cohort, while keeping historical label-collection latency outside the experiment. Delays will remain fixed in observation units when the analysis cohort size changes \cite{souza,alencar}.

The estimand will be the expected error risk of the distributions generating a named cohort $W_s$, conditional on the frozen classifier $f$. It will differ from the realized error count of that finite observed cohort, which becomes exactly known when all its labels arrive. In the primary independent design, reference observations will be sampled independently of training and test observations from a fixed pre-change population $P_0$. The proposed target will be
\begin{align}
\theta_s(f)&=\frac{1}{|W_s|}\sum_{i\in W_s}\mathbb{E}_{(X,Y)\sim P_i}[\ell(f(X),Y)],\\
\theta_0(f)&=\mathbb{E}_{(X,Y)\sim P_0}[\ell(f(X),Y)],\\
\Delta_s&=\theta_s-\theta_0,\qquad H_s=\ind\{\Delta_s>\delta\},
\quad \ell(f(x),y)=\ind\{f(x)\ne y\}.
\end{align}
The primary design margin will be $\delta=0.05$, meaning five additional expected errors per hundred predictions. It will define the scientific benchmark's material degradation threshold rather than a universal deployment tolerance. Margins 0.02 and 0.10 will be prespecified sensitivities, with all conclusions identified by margin. Decision coverage will depend on the distance of excess risk from this boundary relative to interval width, rather than on the margin's magnitude alone; the margin will not be equated with a macro-F1 change \cite{souza,hoeffding}.

Here ``harm'' will mean excess classification error only. A below-margin verdict will mean $\Delta_s\leq\delta$; it will not establish satisfactory absolute accuracy, equitable subgroup performance, safety or business value. Absolute reference and cohort error will therefore accompany every excess-risk display. The benchmark's equal weighting of mechanisms will define an artificial evaluation population, not an estimate of change prevalence in an industry deployment \cite{nist,souza}.

The unit of analysis will be a named prediction cohort, while the independent replication unit will be a final synthetic seed. Every monitoring record will distinguish the input observation time, cohort completion time, label-release time and assessment time; these four clocks will not be treated as interchangeable. At an assessment $q\geq e_s$, displayed cohort-evidence age will mean $q-e_s$ observation units since cohort completion; label-release latency will be recorded separately. A locked study specification will record the target population, loss, margin, eligible evidence and intended claims before data generation or case-study replay \cite{souza,nist}.

\subsection*{Step 2 --- Data Collection}
Controlled synthetic streams will provide event and variable truth. Electricity and the NOAA Weather benchmark described in the stream-benchmarking literature will be the planned natural case studies. Each will require an auditable chronological order, a documented label definition and an identifiable data version before inclusion. Covertype will supply a separate constructed multiclass replay; its source order will not be interpreted as deployment time. A data card will record acquisition conditions, checksums, class counts, missingness and partition boundaries. These retrospective datasets will support cross-domain case comparison, not a claim of representative deployment sampling \cite{souza}.

The proposed data contract will assign an immutable prediction identifier, cohort identifier, model and preprocessing versions, input time, prediction and label-release time to each record. Labels will join one-to-one by prediction identifier, never by row position or nearest timestamp. Exact duplicate identifiers will be rejected; equal feature values alone will not identify duplicates. Conflicting labels, unavailable timestamps and invalid schemas will have explicit failure codes. Naturally observed labels will not be described as error-free population truth. Synthetic labels will follow the declared generator; natural label definitions and known annotation limitations will be recorded before analysis \cite{awsjoin,souza}.

Each final synthetic seed will reserve 10,000 observations for model fitting and a disjoint 10,000-observation pre-change risk reference, shared across its six mechanisms and two reference policies. Its two classifiers will fit the same model-training observations. Each mechanism will generate 60,000 test observations; policies compared within that mechanism will receive exactly the same realized inputs and labels. Mechanisms may share exogenous Gaussian innovations and uniform label-noise draws, but their interventions will change the resulting observations. An additional independent 10,000-observation development sample will fit the common synthetic representation and kernel bandwidth once. Separate random-number namespaces will serve model development, score normalization, threshold selection, held-out null validation, final streams and oracle evaluation. Five model-development seeds and 30 final evaluation seeds will be prespecified. The precision planning in Step~9 will qualify conclusions at that fixed budget; no final-test-dependent replication will occur \cite{souza,cerqueira}.

Independent generator draws, inaccessible to every monitoring policy, will estimate expected reference and cohort risks for evaluation truth. Exact risk will be used where analytically available. Otherwise, $K_g$ will count every distinct regime--classifier risk, including reference risks, evaluated within one seed. A union-bound allocation $\alpha_g/K_g$ and accuracy $\epsilon_g=0.005$ will give
\begin{equation}
\Pr\{|\widehat\theta-\theta|>\epsilon_g\}\leq 2e^{-2M\epsilon_g^2},
\qquad M\geq\frac{\log(2K_g/\alpha_g)}{2\epsilon_g^2},\quad
\alpha_g=\frac{0.01}{30}.
\end{equation}
Repeated distributions will share a regime-level risk estimate. Gradual changes will be mixtures of two declared endpoint distributions, so cohort risk will be a known weighted average of endpoint risks. Excess-risk bounds $[\underline\Delta_s^g,\overline\Delta_s^g]$ will include uncertainty in both the cohort and reference estimates. The latent target $H_s$ in Step~1 will remain binary; its evaluator-only certified label will be
\begin{equation}
\widetilde H_s=
\begin{cases}
1,&\underline\Delta_s^g>\delta,\\
0,&\overline\Delta_s^g\leq\delta,\\
?,&\text{otherwise}.
\end{cases}
\end{equation}
Thus a boundary-uncertain evaluation label will not be confused with a monitoring policy's unresolved verdict. The union allocation will bound Monte Carlo failure across the prespecified 30-seed independent experiment by 0.01, including all declared regime--classifier risks and the three harm margins. No monitoring policy will observe $\widetilde H_s$, use it to determine eligibility or skip a difficult cohort. Oracle failure and sampling uncertainty will be accounted for separately in Step~9. Mechanism parameters will not be adjusted using final-test harm outcomes \cite{hoeffding}.

The collection register will separate source observations from generated evaluation draws and identify which partitions can influence fitting, calibration or final scoring. Shared regime estimates will retain their algebraic dependence: when cohort and risk reference have exactly the same generating law, their expected-risk difference is exactly zero. For the baseline-to-post mixture cohort, $\Delta_s=\overline w_s(\theta_{\mathrm{post}}-\theta_0)$, where $\overline w_s$ is the mean declared mixture weight. Oracle bounds will apply this same linear combination, rather than adding avoidable uncertainty to a term that cancels. Every source and generated partition will have recorded membership and a checksum or reproducible generation specification \cite{hoeffding,souza}.

\subsection*{Step 3 --- Data Preprocessing}
For the synthetic study, imputation, numeric scaling and the kernel representation will be fitted on the independent representation-development sample and frozen across all final seeds. Numeric missing values will use fitted medians; scaled coordinates will be $(x_j-\operatorname{median}_j)/\operatorname{IQR}_j$, with a unit denominator when the fitted interquartile range is zero. Categories will use a fixed one-hot vocabulary with explicit missing and other-category levels. Model configurations will be selected on the five disjoint development seeds: each classifier will fit the first 80\% of its training partition and validate on the final 20\%. The selected configuration will then fit each final seed's complete model-training partition. Shared preprocessing will therefore be independent of final model selection, risk-reference labels and test outcomes. For each natural dataset, candidate preprocessing will fit only the fitting portion of its chronological validation split; after selection, it will be refitted on the full training prefix. Predefined missingness flags will be retained, and encoded contributions will be mapped back to original variables \cite{souza}.

An entirely missing numeric fitting column will receive a fixed zero imputation and unit scale, with its missingness indicator retained; no test statistic will supply its fitting value. Nonfinite numeric values and schema/type violations will be recorded as data-contract failures rather than silently coerced into ordinary measurements. No outlier clipping, target-derived feature, future-based interpolation or outcome-dependent row deletion will be used. Repeated records, temporal aggregation and source identifiers will be audited before partitioning to prevent an artificial chronological split from concealing duplicated observations or future information \cite{souza,nist}.

The primary cohort will contain $b=b_0=1{,}000$ consecutive observations. Cohort sizes 500 and 2,000 will be sensitivities; every analysis will use nonoverlapping cohorts with stride $b$. Both input-reference policies will use a fixed size of 1,000 observations: the final 1,000 rows of the held-out reference prefix, or the immediately preceding 1,000 stream observations. The first rolling reference will therefore also come from that prefix. All observations used for a rolling reference will precede the current cohort. The held-out 10,000-observation predictive-risk reference will remain fixed for both policies. Incomplete cohorts, missing labels and exclusion reasons will be recorded \cite{greco,lukats,souza}.

Cohort membership will be fixed as $W_s=\{(s-1)b+1,\ldots,sb\}$, with membership defined before looking at losses or alarms. A trailing partial cohort will be excluded identically for every policy and listed with its observation count; the core 60,000-observation streams are divisible by every planned cohort size. The preprocessing specification will retain the fitted values, category vocabulary, original-to-encoded feature map and reference-row identifiers, making both transformation leakage and accidental reference overlap auditable \cite{souza}.

\subsection*{Step 4 --- Feature Engineering / Representation}
Original variables will form the principal explanatory representation. Numeric marginal change will use KS distance; Wasserstein distance normalized by the fitted interquartile range will be a comparator. Categorical variables will use Jensen--Shannon divergence of empirical frequencies on the fixed category vocabulary, with base-2 logarithms and $0\log0=0$. A training interquartile range of zero will trigger a declared unit-scale fallback rather than division by zero. The kernel representation will concatenate scaled numeric variables, category indicators and declared missingness indicators \cite{greco,szucs}.
\begin{equation}
D_{j,s}^{\mathrm{KS}}=\sup_x|\widehat F_{j,\mathrm{ref}}(x)-\widehat F_{j,s}(x)|.
\end{equation}
A joint Gaussian-kernel MMD statistic will provide evidence about changes that marginal tests can miss. Original variables will be ranked by the normalized marginal scores defined in Step~5, with ties broken by a fixed variable order. Prediction frequencies, confidence summaries and missingness will be descriptive channels; predicted labels will not substitute for observed correctness. A global joint alarm will not be claimed to identify a changed feature pair. Feature ranking will concern changed marginal distributions, not causal root causes \cite{greco,szucs,driftvis,gretton}.

The monitored object will be the frozen representation $\phi(X)$ rather than an implicitly lossless copy of every raw measurement. Imputation and the other-category mapping can merge distinct raw values. Consequently, equality or non-detection after encoding will not establish raw-data equality. Synthetic inputs are complete numeric variables with an invertible fixed affine transform, so their stated marginal and joint truth will remain observable in that representation. Natural-data missingness and unknown-category rates will be displayed separately from imputed-value distances; marginal rankings there will be descriptive rather than certified raw-variable explanations \cite{gretton,szucs,souza}.

For input-reference policy $r$, let $P_s^X$ be the pooled input-generating distribution of cohort $s$, and $Q_{s,r}^X$ the pooled generating distribution of its reference observations. The native input target will be $D_{s,r}^X=\ind\{P_s^X\ne Q_{s,r}^X\}$. Synthetic marginal truth will be $G_{s,r}=\{j:P_{s,j}^X\ne Q_{s,r,j}^X\}$. Generator-defined mixture weights will determine these distributions, rather than comparing two noisy empirical samples to manufacture evaluation truth. A rolling reference that has assimilated a regime and a fixed reference observed after a return to baseline can legitimately have different targets. Both will retain the same predictive-harm target $H_s$ \cite{window,cerqueira}.

The comparator distances will also have fixed definitions. With fitted scale $a_j=\operatorname{IQR}_j$ or the declared unit fallback, the empirical CDFs will use imputed values in their original numeric units before scaling. Normalized Wasserstein distance will be $D^W_{j,s}=a_j^{-1}\int|\widehat F_{j,\mathrm{ref}}(x)-\widehat F_{j,s}(x)|\,dx$. For categorical probability vectors $p,q$, Jensen--Shannon divergence will be $\operatorname{JS}(p,q)=\{\operatorname{KL}_2(p\Vert m)+\operatorname{KL}_2(q\Vert m)\}/2$, $m=(p+q)/2$, omitting zero-probability summands. The feature dictionary will distinguish original variables, encoded coordinates and descriptive missingness rates; encoded coordinates will not be counted as independent ground-truth variables \cite{szucs,gretton}.

\subsection*{Step 5 --- Proposed Model / Method}
The monitor will maintain input-change evidence separately from predictive-harm evidence. For reference observations $u_1,\ldots,u_m$ and cohort observations $v_1,\ldots,v_n$, the proposed joint statistic will use the usual squared-MMD U-statistic form \cite{gretton}:
\begin{align}
\widehat{\mathrm{MMD}}_u^2={}&\frac{\sum_{i\ne j}^{m} k(u_i,u_j)}{m(m-1)}
+\frac{\sum_{i\ne j}^{n} k(v_i,v_j)}{n(n-1)}\\
&-\frac{2\sum_{i=1}^{m}\sum_{j=1}^{n}k(u_i,v_j)}{mn},\nonumber\\
k(u,v)&=\exp\!\left(-\frac{\|u-v\|^2}{2\sigma^2}\right).
\end{align}
Bandwidth will be the median nonzero pairwise distance on a reproducibly selected 500-row subset of the independent representation-development sample and will then be frozen. For a natural dataset or Covertype replay, the corresponding subset will come only from its model-training prefix. If no positive distance exists, unit bandwidth will be used. To bound computation, MMD will use at most 200 reproducibly selected observations from each input window; the same relative-window indices will be shared across policies. Marginal statistics will use all observations. Negative U-statistic values will be retained during calibration, rather than silently clipped \cite{gretton}.

Unbiasedness for a population squared MMD will be invoked only when observations are i.i.d. within each of the two independent samples. Transition cohorts can be independent without being identically distributed, and dependent controls violate independence; in those settings the same formula will be a descriptive score, not an automatically unbiased estimate with the original test guarantees. The declared null calibration will determine its alarm operating point \cite{gretton}.

Component normalization will use the independent score-normalization bank $\mathcal C$. For each component $j$, including joint score $J$, define its median $c_j$, interquartile range $q_j$, and positive scale
\begin{align}
c_j&=\operatorname{median}_{D\in\mathcal C_j}D,\\
v_j&=\begin{cases}q_j,&q_j>0,\\
\max\{10^{-6},\max_{D\in\mathcal C_j}|D-c_j|\},&q_j=0,
\end{cases}\\
Z_j(D)&=(D-c_j)/v_j,\\
S_s&=\max\{Z_1(D_{1,s}),\ldots,Z_p(D_{p,s}),Z_J(J_s)\},\qquad
A_s=\ind\{S_s>\tau\}.
\end{align}
The entire maximum will receive a separately selected threshold. These dimensionless component scores will not be claimed to be p-values. Unlike an empirical-CDF transformation, the affine normalization will not clip all scores above the observed calibration maximum to one; extreme new observations will remain distinguishable. The constant-score fallback will be fixed before testing, and ties in feature ranking will follow the declared variable order \cite{cerqueira,gretton}.
The proposed null banks will contain 2,000 score-normalization windows, 5,000 threshold-selection windows and 2,000 held-out validation windows per reference policy and null process. These will comprise, respectively, 20, 50 and 20 independent 100-window stationary trajectories, each with its own pre-change input reference. Within-trajectory dependence and the reference-update schedule will be preserved. The threshold will be the smallest empirical value with no more than 1\% of threshold-selection scores strictly above it; ties will make the calibration conservative. Validation uncertainty will resample whole trajectories rather than treat shared-reference windows as independent. Ten raw alerts per 1,000 windows will be a target averaged over reference draws, with achieved rate, uncertainty, per-trajectory burden and calibration cost reported. It will not guarantee a 1\% conditional rate for every particular frozen reference. Validation will not be used to retune a threshold, and these rates will not be called lifetime false-alarm guarantees \cite{lukats,cerqueira}.

Because the synthetic representation and null laws will be fixed, these banks will be generated once per reference policy, cohort size and null process and reused across classifiers, label schedules and final seeds. Window-size sensitivities and detector ablations will receive their own score normalization and threshold selection from the appropriate development bank; they will not reuse a threshold calibrated for a different statistic or sample size. Inference will be conditional on this frozen calibration design. Null observations will be generated sequentially, retaining statistics and seeds rather than full trajectories. All detector families will use the same source trajectories, and label-delay variants will reuse input scores. These choices will limit computation without treating reused calibration data as extra test replication \cite{souza,cerqueira}.

For fully released cohorts, empirical excess loss will be
\begin{equation}
\widehat\Delta_s=\frac{1}{|W_s|}\sum_{i\in W_s}\ell(f(x_i),y_i)
-\frac{1}{n_0}\sum_{i\in W_0}\ell(f(x_i),y_i),\qquad n_0=10{,}000.
\end{equation}
In the primary independent synthetic design, a pointwise Hoeffding interval will be applied directly to the weighted difference of the two means. Conditional on $f$, independent bounded losses need not have equal expectations within a transition cohort. The squared ranges of the weights sum to $1/|W_s|+1/n_0$. For $\alpha=0.05$, the proposed interval will be
\begin{align}
a_s&=\sqrt{\frac{\log(2/\alpha)}{2}
\left(\frac{1}{|W_s|}+\frac{1}{n_0}\right)},\\
[L_s,U_s]&=[\widehat\Delta_s-a_s,\widehat\Delta_s+a_s]\cap[-1,1].
\end{align}
At $|W_s|=1{,}000$ and $n_0=10{,}000$, the planned half-width will be approximately 0.045. The rule will therefore require empirical excess loss above approximately 0.095 for a primary harm declaration, or at most approximately 0.005 for a below-margin declaration; this is a conservative discrimination band around $\delta=0.05$, not evidence of high power near that boundary. These are analytical design calculations, not experimental results. The rule will return harm-supported if $L_s>\delta$, below-margin if $U_s\leq\delta$, and unresolved otherwise. Eligibility will require every label in a complete declared cohort. Levels 90\% and 99\% will be sensitivities. Each decision will use a predeclared checkpoint; these pointwise intervals will not be advertised as anytime-valid certificates or as 95\%-accurate statements conditional on having made a decision \cite{hoeffding,selective}.

To make evidence eligibility explicit, let $c_s(q)=|W_s|^{-1}\sum_{i\in W_s}\ind\{r_i\leq q\}$ be released-label coverage. Let $\nu_s(q)\in\{0,1\}$ indicate record and one-to-one join validity. The monitor will require $c_s(q)=1$ and $\nu_s(q)=1$, then apply the interval rule:
\begin{equation}
B_s(q)=\begin{cases}
\bot,&c_s(q)<1\text{ or }\nu_s(q)=0,\\
1,&c_s(q)=1,\ \nu_s(q)=1,\ L_s>\delta,\\
0,&c_s(q)=1,\ \nu_s(q)=1,\ U_s\leq\delta,\\
\bot,&\text{otherwise}.
\end{cases}
\end{equation}
An invalid record will raise a pipeline-failure status and invoke the prespecified run-repair procedure; it will never be relabelled below-margin or silently omitted from a completed comparison. This rule will produce a cohort evidence record containing the input alarm, reference identity, label coverage, interval, verdict and evidence age \cite{awsjoin,alencar,nist}.

This interval also supplies an analytical reference for interpretation. Under the independent bounded-loss assumptions, for a fixed cohort with $H_s=0$, a wrong positive declaration requires $\widehat\Delta_s-\Delta_s>a_s$, and hence has probability at most $\alpha/2=0.025$. Adding the condition $A_s=0$ cannot increase that unconditional probability. Linearity of expectation preserves this bound for the all-cohort mean wrong-added-declaration frequency, despite reuse of the historical reference. It is not a bound on the fraction wrong among selected positive declarations, the certified subset, or the chance of any error over the entire stream. Empirical error analysis will investigate operating behaviour and assumption failures rather than claim this existing concentration bound as algorithmic novelty \cite{hoeffding}.

Dependent stationary controls will use a separate reference trajectory and an exploratory two-sample moving-block bootstrap with 2,000 replicates. Within each loss sequence, starts will be sampled uniformly from all complete overlapping blocks of length $l$; sampled blocks will be concatenated and truncated to the original sequence length. Reference and cohort resampling will be independent, preserving their separate means and avoiding artificial mixing of different regimes. The 2.5th and 97.5th percentiles of replicate mean differences will define the interval. Development controls will select $l\in\{5,10,20,50\}$ by proximity of empirical coverage to 0.95, breaking ties toward shorter blocks. A separate validation bank will use 2,000 independent trajectory pairs, each with a new reference and cohort; reusing one reference across these pairs would not supply independent coverage trials. Stationary excess-risk truth is exactly zero because the reference and cohort have the same marginal law. Validation will report binomial coverage uncertainty and interval width, without selecting another block length. The 93\%--97\% empirical coverage band will be an exploratory diagnostic, not proof of nominal coverage. Failure will be reported rather than repaired using final outcomes; stationary validation will not establish validity inside nonstationary transitions \cite{kunsch}.

For natural cases, the same separate-segment resampling will provide exploratory uncertainty only, conditional on local stability within each segment and negligible dependence between the historical reference and current cohort. Cross-segment dependence will not be claimed to have been modelled, and resampling will never concatenate the two regimes into a supposed stationary sequence. Population harm will remain unverified; neither these intervals nor the dependent controls will inherit the independent-design Hoeffding guarantee \cite{kunsch,souza}.

\subsection*{Step 6 --- Baseline Methods}
Distributional detector comparators will be marginal KS, normalized Wasserstein and joint MMD families. All will receive independently selected thresholds under the same null-rate target. Neither DriftLens nor proactive adaptation will be claimed as exactly reproduced in a frozen tabular-classifier study. Their literature findings motivate the comparison but concern different tasks \cite{greco,szucs,guerrero}.

PUDD will inform the released-label channel because its PU score $u_i=1-f_{y_i}(x_i)$ depends on the observed true class. It will not be redefined as an unlabeled confidence score. The core comparison will isolate verdict gating using the simple detector families above; it will not claim superiority to PUDD, conformal model reuse or the fourteen-detector KDD benchmark without a separate faithful reproduction. Their original adaptive-learning scores will not be imported as comparable harm-verdict results \cite{pudd,conformal,cerqueira}.

Let $B_s$ be the eligible interval verdict from Step~5, taking values 1, 0 or $\bot$ for harm-supported, below-margin or unresolved. The primary comparator will gate positive harm declarations only. A below-margin verdict will remain available in both policies; an ineligible or inconclusive cohort will remain unresolved in both:
\begin{equation}
V_s^{\mathrm{separated}}=B_s,\qquad
V_s^{\mathrm{gated}}=
\begin{cases}\bot,& B_s=1\text{ and }A_s=0,\\ B_s,&\text{otherwise}.\end{cases}
\end{equation}
Keeping below-margin verdicts identical will focus the matched comparison on harm recognition. Its coverage ordering will remain known by construction; the empirical question will concern the correctness and mechanism-specific frequency of additional harm declarations. This comparator will be a declared policy experiment, not a claim that all deployed monitoring systems implement such a gate or that gating saves label-acquisition cost \cite{alencar,solozobov}.

The primary question will be estimation-based: how often are these additional harm declarations correct or wrong, especially after change onset? No test will treat the known nonnegative coverage ordering as an empirical discovery. The proposed comparison will succeed as an evaluation if it produces auditable all-cohort bounds, uncertainty and mechanism-specific failure evidence; a claim of practical improvement will additionally require an application-defined benefit/error tolerance that this general benchmark does not possess. A null or inconclusive conclusion will be legitimate if additional declarations are absent, wrong-decision uncertainty is substantial, or planned precision is not attained. No application-independent monetary value will be assigned to this trade-off \cite{window,pudd,selective,nist}.

Secondary verdict comparators will include the input-only proxy $V_s=A_s$, an eligible point-loss rule $\ind\{\widehat\Delta_s>\delta\}$, and a point-loss rule requiring $A_s=1$ for a harmful verdict but assigning below-margin to other eligible cases. A stricter all-verdict gate, $V_s^{\mathrm{strict}}=B_s$ when $A_s=1$ and $\bot$ otherwise, will separately show the effect of withholding below-margin verdicts as well. All label-based rules will defer when ineligible. Always-below-margin and always-unresolved controls will expose harm prevalence and total-deferral behaviour. Every policy will receive the same classifier, cohort, label schedule and checkpoint; input-only output will be explicitly described as a harm proxy, not the intended native task of a distributional detector \cite{alencar,solozobov}.

The loss-only interval control will use precisely the same harm rule as the proposed evidence-separated display. Their verdicts will be required to be identical; this will be an identity check rather than a statistical superiority hypothesis. Input evidence will add context without changing that rule. The point-loss and point-gated comparisons will be secondary ablations, so changes in uncertainty handling will not explain the primary gating contrast \cite{solozobov}.

Baseline roles will remain distinct. The marginal and joint distance families will test the input channel under matched stationary alert targets; the two primary verdict policies will differ only in the positive-alarm gate applied to the same $B_s$. Point-loss variants will test the uncertainty component, and constant-verdict controls will expose prevalence and deferral. All verdict policies will read one common cohort evidence record at the same checkpoint. A shared null-alert target will not be described as equal detection power, equal computation or individual-feature error control \cite{cerqueira,selective}.

\subsection*{Step 7 --- Experimental Design}
The primary design will cross two input references with delays $d\in\{0,2,5\}$ and two frozen classifiers on identical controlled streams. A logistic model will select regularization from $\{0.1,1,10\}$; the tree will select depth from $\{3,5\}$ on the five development seeds. The fitting objective will be mean chronological validation error across those seeds, with ties resolved toward stronger regularization or shallower depth. Configurations will then be frozen for all final seeds. The primary synthetic data will use ten numeric variables and observations independent across time; autoregressive controls will form a separate robustness tier \cite{souza,cerqueira}.

Before final evaluation, a dated protocol manifest will freeze the development/final seed lists, dataset hashes, split indices, event schedules, parameter grids, solver settings, calibration banks, thresholds, metric denominators and analysis scripts. No seed, mechanism, classifier or cohort will be removed for an unfavourable result. A failed numerical or data-contract check will suspend the affected run, retain its failure log and require an outcome-blind correction followed by rerunning the same complete paired seed. Checkpoint records will be append-only; corrected or later labels will not retrospectively replace the recorded evidence available at an earlier checkpoint \cite{nist,awsjoin,souza}.

The core design will contain six mechanisms, two frozen classifiers and two input references: 24 accuracy conditions per seed at the primary $h=2,d=0$ checkpoint. Three label schedules will produce 72 scheduled condition cells per seed, but will reuse the same six realized input/label streams. Thirty final seeds therefore imply 180 mechanism-stream realizations and 10.8 million test observations, not 2,160 independent experiments. For the primary gradual transition, the exact schedule will be $w_i=(i-20{,}001)/4{,}999$ for $20{,}001\leq i\leq25{,}000$, with weights zero before and one afterward. A correction to a shared routine will invalidate and rerun every affected completed condition using its original seed; it will not be applied only to an unfavourable case \cite{souza,cerqueira}.

\begin{longtable}{L{4.2cm}L{11.7cm}}
\caption{Planned controlled conditions and evaluation truth.}\label{tab:conditions}\\
\toprule\rowcolor{headgray}\textbf{Condition}&\textbf{Purpose and truth}\\\midrule\endfirsthead
\toprule\rowcolor{headgray}\textbf{Condition}&\textbf{Purpose and truth}\\\midrule\endhead
\bottomrule\endlastfoot
Stationary & Measure raw null alerts on independent and autoregressive stationary controls.\\[3pt]
Irrelevant marginal change & Shift an independently generated nuisance variable; determine actual classifier harm from independent risk truth rather than its nuisance designation.\\[3pt]
Relevant marginal change & Shift a predictive coordinate while retaining the label rule; include abrupt and gradual transitions.\\[3pt]
Recurring correlated change & Record onset, completion and return, separating directly manipulated coordinates from all changed marginals.\\[3pt]
Dependence-only change & Vary correlation with invariant univariate marginals; score joint detection without inventing feature-pair localization output.\\[3pt]
Conditional-only change & Alter the label mechanism with invariant input distribution; assess delayed harm verdicts without expecting an input alarm.\\
\end{longtable}

Independent $Z_1,Z_2\sim\mathcal N(0,1)$ will permit the dependence-only construction
\begin{equation}
X_1=Z_1,\qquad X_2=\rho_sZ_1+\sqrt{1-\rho_s^2}\,Z_2,\qquad |\rho_s|<1.
\end{equation}
Both marginals will stay standard normal. Baseline correlation will be $\rho=0.3$, and the binary label rule will be $Y\sim\operatorname{Bernoulli}(\operatorname{sigmoid}(X_1+X_2))$; the other eight coordinates will be independent standard normals. The primary marginal shift will be one standard deviation, applied to coordinate 10 for nuisance change or coordinate 1 for predictive change, retaining the same $P(Y\mid X)$. Recurrence will shift coordinates 1 and 2 by 1 and 0.5 respectively, then return to baseline. Dependence-only change will set $\rho=-0.6$, while conditional-only change will negate the label-rule coefficients and preserve $X$. Amplitudes 0.5 and 2 and correlations 0 and 0.6 will be prespecified secondary sensitivities, not selected using final harm outcomes. The actual harm state will always come from independent risk truth \cite{driftvis,gretton}.

Abrupt onset will be at test observation 20,001. The primary predictive marginal-change condition will use a gradual mixture through observations 20,001--25,000, $P_i=(1-w_i)P_{\mathrm{pre}}+w_iP_{\mathrm{post}}$, with $w_i$ increasing linearly from zero to one. The nuisance, dependence-only and conditional-only conditions will use abrupt onset; recurrence will return at observation 40,001. An abrupt predictive-change variant will be secondary. The same observation-index events will apply at every cohort size. A dependent stationary control will use $X_i=aX_{i-1}+\varepsilon_i$ with $a=0.5$. Its covariance $\Sigma_0$ will have unit diagonal, entries $(1,2)$ and $(2,1)$ equal to 0.3, and all other off-diagonal entries zero. Initialization $X_0\sim\mathcal N(0,\Sigma_0)$ and independent innovations $\varepsilon_i\sim\mathcal N(0,(1-a^2)\Sigma_0)$ will give the declared stationary marginal law from the first observation; label-noise uniforms will be independent over time. Independent stationary controls will use $a=0$. Generator event truth will remain distinct from classifier harm \cite{souza,cerqueira}.

For cohort completion $e_s$, the primary assessment time will be $q_s(h)=e_s+hb_0$ with $h=2$; $h=0$ and $h=5$ will be sensitivities. A verdict will use the cohort's own stored input alarm $A_s$, not an alarm from a later window, and only labels released by $q_s(h)$. With complete-label eligibility and fixed delays, $d>h$ structurally prevents an eligible loss-based verdict. This will be reported as evidence unavailability rather than an estimated accuracy loss. Later labels will not rewrite earlier checkpoint results. After the last input, the clock will continue through the maximum planned checkpoint without generating extra test observations. Permanent missingness or selective label acquisition will not be introduced into this complete-release experiment \cite{alencar}.

With a constant delay and complete-cohort eligibility, the first eligible time will be exactly $e_s+db_0$; the loss estimate will then be identical at all later checkpoints. This is a controlled availability experiment, not a test of partial-evidence accumulation or a delay-induced change in the classifier. Random or feature-dependent delays would require a separate extension. The results will distinguish this structural label-waiting time from additional deferral caused by an interval crossing the harm margin or by alarm gating \cite{pudd,alencar,window}.

Each natural case study will preserve record order, using the first 30\% of observations for fitting, the next 20\% for reference and calibration, and the last 50\% for testing. The two historical prefixes will be treated as fully labelled archives before the retrospective test clock starts; the artificial release schedule will then restrict test labels only. This will not simulate calendar-time acquisition of the historical labels. The source will be required to provide at least five complete test cohorts and both classes in the training prefix; any exclusion will be recorded before outcome analysis. Model selection will use an 80/20 chronological split inside the fitting prefix and the same configuration grids and tie rules as the synthetic design, followed by refitting on the complete fitting prefix. No source-specific thresholds or model choices will be tuned on test data. Score normalization and threshold selection will use successive disjoint halves of the reference prefix. Because these halves may provide too few original windows for a rare-tail threshold, 2,000 and 5,000 moving-block pseudo-windows will respectively form exploratory reference-based calibration banks. Within each half, contiguous blocks of length 20 will be sampled with replacement to form pseudo-trajectories containing both the input reference and current window; the declared fixed or rolling update rule will then be applied. This block length will be an exploratory convention, not a data-driven optimum. Bootstrap blocks will never cross the half-partition boundary. They will not count as new independent observations or certify a stationary 1\% false-alert rate. Risk estimation will use the entire held-out labelled prefix, with its actual size replacing $n_0$ in Step~5 \cite{souza,kunsch}.

Natural alert frequency will be labelled alert burden, and realized losses will be observable after release, but population-harm truth will not be treated as known. Synthetic $I$, FR and FH scores will therefore not be assigned to natural data. Electricity and Weather results will be shown separately, without pooling them as exchangeable deployments. Covertype will receive one recorded seeded permutation, followed by disjoint 30/20/50 partitions; this replay will illustrate multiclass monitoring under an artificial ordering, without certified temporal events or population-harm truth. Multinomial logistic regression and a multiclass tree will use the same selection grids, with all source classes required in training. Binary logistic predictions will use a 0.5 cutoff, multiclass predictions will use the largest probability, and class ties will follow a fixed class order. Logistic regularization will mean an L2 penalty $\lambda\|\beta\|_2^2/2$ added to mean log loss; the intercept will be unpenalized. Trees will use Gini impurity, a minimum leaf size of 20 and the selected depth. Transfer claims will be limited to the tested domains and representations \cite{souza}.

\subsection*{Step 8 --- Evaluation Metrics}
Let $N_{\mathrm{all}}$ denote all complete declared test cohorts, whether or not their oracle truth is certified. The latent all-cohort primary quantities can be written directly. With $Q_s=\ind\{A_s=0,B_s=1\}$,
\begin{equation}
g_{\mathrm{correct,all}}=\frac{\sum_s Q_sH_s}{N_{\mathrm{all}}},\qquad
g_{\mathrm{wrong,all}}=\frac{\sum_s Q_s(1-H_s)}{N_{\mathrm{all}}}.
\end{equation}
Their sum is the observable added-declaration frequency; the oracle bounds address only how that frequency divides into correct and wrong cases. Every result table will identify its denominator, eligibility rule, truth source and uncertainty type. Identification width, label-unavailable fraction and monitor deferral will be shown separately, so absent truth and absent monitoring evidence cannot improve an apparent accuracy score \cite{hoeffding,selective}.

The primary outcome will be the all-cohort correct and incorrect additional-declaration frequencies, reported as identification bounds when truth is uncertain, together with all-cohort coverage and decision availability. Certified-subset scores will be secondary diagnostics because selecting only truth far enough from the margin can change the difficulty and harm prevalence of the evaluated population. Let $\mathcal E=\{s:\widetilde H_s\in\{0,1\}\}$ and $N=|\mathcal E|$. On this common diagnostic set for every policy, let $D$ count decisive verdicts and $E$ count verdicts differing from $\widetilde H_s$. The study will report
\begin{equation}
I=\frac{E}{N},\qquad C=\frac{D}{N},\qquad U=1-C,\qquad
E_{\mathrm{conditional}}=\frac{E}{D}\quad(D>0).
\end{equation}
The always-unresolved rule will have $I=0$, $C=0$ and undefined conditional error. Low incorrect-claim frequency alone will therefore not establish an improvement. Conditional error will be compared only at genuinely shared operating points; deterministic proxy output will not be assigned a fictitious risk--coverage curve \cite{selective}.

For the matched positive-gating contrast, the additional-declaration set will be $\mathcal Q=\{s\in\mathcal E:A_s=0,\ B_s=1\}$. Its correct and incorrect counts, conditional error and distribution across mechanisms will be reported. The certified-subset diagnostic components will be
\begin{align}
g_{\mathrm{correct}}&=(C_{\mathrm{separated}}-I_{\mathrm{separated}})
 -(C_{\mathrm{gated}}-I_{\mathrm{gated}}),\\
g_{\mathrm{wrong}}&=I_{\mathrm{separated}}-I_{\mathrm{gated}},\qquad
C_{\mathrm{separated}}-C_{\mathrm{gated}}=g_{\mathrm{correct}}+g_{\mathrm{wrong}}.
\end{align}
Equivalently, $g_{\mathrm{correct}}=|\{s\in\mathcal Q:\widetilde H_s=1\}|/N$ and $g_{\mathrm{wrong}}=|\{s\in\mathcal Q:\widetilde H_s=0\}|/N$. Both will be nonnegative by construction. The study will estimate their magnitude rather than test this ordering as a discovery. Conditional error on $\mathcal Q$ will be undefined when no additional declarations occur. These diagnostic quantities will not gain credit for additional below-margin decisions, which occur only in the secondary strict-gate comparison. A coverage increase alone will not establish practical benefit \cite{alencar,selective}.

Both all-cohort bounds and certified-subset diagnostics will have mandatory pre-onset/post-onset breakdowns; certified harmful/below-margin strata will additionally show their counts and uncertain-truth frequency. Stationary controls will use the observation-20,001 reporting cut without inventing a drift event. These summaries will distinguish recovered post-change harm declarations from incorrect pre-change declarations. No favourable mechanism selection or post-hoc prevalence weighting will replace the prespecified primary summary \cite{window,cerqueira}.

Boundary-uncertain truth will not silently remove difficult cases. Let $N_{\mathrm{all}}$ include all complete cohorts, $E_{\mathrm{known}}$ count errors on decisive truth, and $D_{\mathrm{uncertain}}$ count decisions on boundary-uncertain truth. All-cohort coverage will use $D_{\mathrm{all}}/N_{\mathrm{all}}$, and the logically possible all-cohort error frequency will be bounded by
\begin{equation}
\frac{E_{\mathrm{known}}}{N_{\mathrm{all}}}\ \leq\ I_{\mathrm{all}}\ \leq\
\frac{E_{\mathrm{known}}+D_{\mathrm{uncertain}}}{N_{\mathrm{all}}}.
\end{equation}
These are identification bounds conditional on the simultaneous oracle-coverage event, whose failure probability is at most 0.01 under the stated generator design. They are not unconditional logical guarantees or sampling confidence intervals. Counts and truth-bound widths will accompany the certified subset; the monitoring policy will still act on every complete cohort, including those whose evaluation truth is uncertain. No mechanism will be excluded because it produces an inconvenient harm prevalence \cite{hoeffding,selective}.

The same accounting will accompany the primary added-declaration components. Let $q_1$, $q_0$ and $q_?$ count all additional declarations with certified harmful, certified below-margin and boundary-uncertain evaluation labels. On all complete cohorts,
\begin{equation}
g_{\mathrm{correct,all}}\in\frac{[q_1,q_1+q_?]}{N_{\mathrm{all}}},\qquad
g_{\mathrm{wrong,all}}\in\frac{[q_0,q_0+q_?]}{N_{\mathrm{all}}}.
\end{equation}
Their sum will equal $(q_1+q_0+q_?)/N_{\mathrm{all}}$, so the two upper bounds cannot generally hold simultaneously. These all-cohort quantities will be the primary estimands. Each of the four endpoints will be retained per seed and condition; oracle-uncertain cohorts will keep the same denominator. The certified-subset estimates above will not be substituted for them. Monitoring deferral and uncertain evaluation truth will remain separate missing-information mechanisms \cite{hoeffding}.

False reassurance and unnecessary harm declarations will have explicit truth-based denominators:
\begin{equation}
\mathrm{FR}=\frac{\sum_{s\in\mathcal E}\ind\{\widetilde H_s=1,V_s=0\}}{\sum_{s\in\mathcal E}\ind\{\widetilde H_s=1\}},\qquad
\mathrm{FH}=\frac{\sum_{s\in\mathcal E}\ind\{\widetilde H_s=0,V_s=1\}}{\sum_{s\in\mathcal E}\ind\{\widetilde H_s=0\}}.
\end{equation}
Because the primary policies retain identical below-margin verdicts, their false-reassurance counts and FR values will be identical whenever defined. This will be an identity check, not a claimed improvement; proxy and other secondary policies can differ because they map evidence to verdicts differently \cite{selective}.
Here unresolved is neither 0 nor 1; deferral will also be reported separately within each harm stratum. A zero denominator will produce an undefined metric with its count displayed. Time to complete eligible evidence will follow the known release schedule. For policy $p$, define $t_s^{p,\star}$ as the first decisive time among the checkpoints $h\in\{0,2,5\}$, measured after cohort completion, and $t_s^{p,\star}=\infty$ if none is decisive. With $H=5{,}000$, the restricted time without a decision and no-decision fraction will be
\begin{equation}
\overline W_H^p=\frac{1}{N_{\mathrm{all}}}\sum_s\min(t_s^{p,\star},H),\qquad
\rho_H^p=\frac{1}{N_{\mathrm{all}}}\sum_s\ind\{t_s^{p,\star}=\infty\}.
\end{equation}
The infinite-time cases will be classified by evidence unavailability, inconclusive interval or positive-alarm gate. Under the complete-release design, $B_s$ and $A_s$ never change after full release, so interval- or gate-induced deferral is permanent for that cohort within this protocol. For either primary label-based policy, the checkpoint grid contains every declared release delay, implying $\overline W_H^p=db_0(1-\rho_H^p)+H\rho_H^p$. Thus this time summary is derived from delay and final decision coverage, not an independent performance endpoint. It will not be presented as eventual decision time or analysed using an assumption that all cases will later decide. Repeated checkpoint descriptions will not imply simultaneous coverage of the monitoring intervals \cite{alencar,selective}.

Native input-state recall and false-positive frequency will first compare $A_s$ with reference-specific truth $D_{s,r}^X$, reporting denominators and undefined cases. This will separate failure to detect a difference from a legitimate absence of difference against that reference. For secondary onset-event summaries, matching will use the first unused alarm after an onset that creates a nonzero native difference, at most once per event. Alarm time will be cohort completion. Abrupt-event tolerance will be 5,000 observations; gradual tolerance will include the transition and 5,000 observations after completion. Additional unmatched alerts will be event-matching false positives, which need not be false alarms about an existing reference difference. Return-to-baseline behaviour will be shown as onset or cessation of native difference, as appropriate to the reference, rather than automatically scored as a missed positive alarm. Eligible-event counts will be shown per policy, and event F1 values with different target sets will not form a common ranking. Tolerances of 1,000 and 3,000 observations will be sensitivities and remain fixed across cohort sizes \cite{window,cerqueira}.
\begin{align}
P_e&=\frac{TP}{TP+FP},&R_e&=\frac{TP}{TP+FN},&F1_e&=\frac{2TP}{2TP+FP+FN},\\
\mathrm{FA}_{1000}&=1000\frac{\text{raw alerts in verified stationary windows}}{\text{verified stationary windows}},
&\mathrm{Delay}_e&=\widehat t_e-t_e.
\end{align}
Detected-input-event delay will be shown with the miss rate and will be undefined when no event is detected. Harm timing will use the cohort-specific decision-time summaries above: a positive verdict will count as correct harm recognition only for an evaluator-certified harmful cohort. Incorrect positive claims will not earn credit as early recognition, and a verdict about an old cohort will not establish the current model's status. A return to baseline will be an input event but will not automatically count as harmful. Conditional-only changes will be excluded from input-event recall and assessed through harm outcomes \cite{driftvis,cerqueira}.

Marginal localization will use the reference-specific set $G_{s,r}$ defined in Step~4 and top-$k$ ranked variables $T_s$, with $k=\min(3,p)$ and average precision as a supplement:
\begin{equation}
P@k=\frac{|T_s\cap G_{s,r}|}{k},\qquad R@k=\frac{|T_s\cap G_{s,r}|}{|G_{s,r}|}.
\end{equation}
Ranking quality will be evaluated on every cohort with nonempty $G_{s,r}$, not only detected cohorts. A rank alone will not constitute a claimed change: the displayed screened-feature set will be $\widehat G_{s,r}=T_s\cap\{j:Z_j(D_{j,s})>\tau\}$. False feature flags per cohort will be $N_{\mathrm{all}}^{-1}\sum_s|\widehat G_{s,r}\setminus G_{s,r}|$, accompanied by the fraction of cohorts with any flag. An MMD-only alarm can therefore have an empty screened-feature set. The maximum-score threshold is calibrated under the complete stationary null; it will not be advertised as strong family-wise error or false-discovery control for individual variables when other variables have changed. Empty marginal truth will have undefined localization recall. Ranking, screening and causal explanation will remain distinct outputs; none will identify a changed feature pair or establish a causal root cause \cite{szucs,causal,gretton}.

Predictive description will include bounded error, class prevalence, per-class recall and fixed-universe macro-F1:
\begin{equation}
\mathrm{MacroF1}=\frac{1}{K}\sum_{c=1}^{K}\frac{2TP_c}{2TP_c+FP_c+FN_c}.
\end{equation}
For a class with zero denominator, its F1 contribution will be defined as zero by a declared reporting convention; unsupported-class counts will accompany the fixed $K$-class average. Runtime and peak memory will include preprocessing, detection, uncertainty calculation and explanation, with calibration cost and steady-state monitoring reported separately. Active label requests, intervention costs and realized retraining benefits will not be outcomes because the core design has fixed label schedules and frozen classifiers \cite{guerrero,souza}.

\subsection*{Step 9 --- Statistical / Comparative Analysis}
Independent synthetic seeds will be replication units, conditional on the common frozen development configuration. A seed will retain its entire vector of shared streams, classifiers, input references and oracle draws. The primary target will average all-cohort added-declaration frequencies equally over six mechanisms, two classifiers and two input-reference policies at $h=2,d=0$. This fixed 24-condition average will characterize the benchmark mixture, not population prevalence in a deployment. All four identification endpoints from Step~8 will be averaged using the same weights. Unlike certified-subset ratios, these quantities remain defined even if no cohort has decisive evaluation truth. Secondary amplitudes, dependent controls, natural cases and window-size variants will not enter this average \cite{souza,hoeffding}.

To answer the reference-policy part of the question, the report will additionally retain each reference's mechanism-specific all-cohort proxy-error and added-declaration bounds. For a quantity bounded by $[l_F,u_F]$ under the fixed reference and $[l_R,u_R]$ under the rolling reference on the same seed, its fixed-minus-rolling contrast will lie in $[l_F-u_R,u_F-l_R]$ on the oracle-coverage event. These conservative contrast bounds will use paired seed resampling for approximate uncertainty; they will not inherit the two-component primary simultaneous guarantee. A pooled average alone will never support a conclusion that one reference is uniformly preferable across mechanisms \cite{window,hoeffding}.

Delay $d=2$ will reproduce the same eligible accuracy contrast and supply only timing information. Delay $d=5,h=2$ will be descriptive because every label-based policy is structurally unresolved. A reference-by-delay accuracy interaction will not be tested as if identifiable. Reused cohorts, overlapping reference observations and repeated checkpoints will never increase the number of independent replications \cite{alencar,souza}.

For seed $j$ and component $k\in\{\mathrm{correct},\mathrm{wrong}\}$, let $G_{jk}\in[0,1]$ denote the latent all-cohort added-declaration frequency and $[\ell_{jk},u_{jk}]$ its condition-averaged identification bounds. The target will be $\mu_k=\mathbb E[G_{jk}\mid\text{frozen development configuration}]$. On the simultaneous oracle-coverage event, $\ell_{jk}\leq G_{jk}\leq u_{jk}$ for every final seed. A conservative finite-sample sampling envelope will use the one-sided bounded-mean inequality:
\begin{align}
\eta_n&=\sqrt{\frac{\log(4/\alpha_s)}{2n}},\qquad n=30,\quad \alpha_s=0.04,\\
\mathcal J_k&=\left[\max(0,\overline\ell_k-\eta_n),\;
\min(1,\overline u_k+\eta_n)\right].
\end{align}
Four one-sided sampling bounds, two per component, will have total failure probability at most 0.04. Combined by a union bound with the prespecified 0.01 oracle budget, both $\mathcal J_k$ will cover their latent seed-population means with probability at least 0.95 under the independent synthetic design. This statement will be conditional on the frozen configuration and correct generator assumptions; it will not transfer to natural or dependent cases. At 30 seeds, $\eta_n\approx0.277$, so this exact guarantee can be too wide to support a practical conclusion. It will remain reported even if a narrower approximate interval appears favourable. The independent-design analytical constraint $\mu_{\mathrm{wrong}}\leq0.025$ from Step~5 will be shown alongside the empirical envelope, clearly labelled assumption-derived rather than measured performance; the difficult remaining empirical quantity will be the yield of correct additional declarations \cite{hoeffding}.

For useful precision description, paired seed-level percentile bootstrap resampling with 10,000 replicates will also report the mean identification interval and uncertainty in its endpoints. A replicate will retain each sampled seed's complete condition vector. The 1st percentile of the mean lower endpoint and 99th percentile of the mean upper endpoint will form an approximate component envelope; four such one-sided endpoint tails will target a nominal 0.04 sampling allocation. Finite-sample bootstrap coverage will not be asserted, particularly with sparse or boundary outcomes. An all-zero bootstrap distribution will be labelled degenerate, not evidence that a population event is impossible; the finite-sample envelope above will still apply. Neither interval will treat uncertain oracle labels as known or replace them by a midpoint \cite{souza,hoeffding}.

Secondary reference comparisons will use the observable disagreement frequency $J_r=N_{\mathrm{all}}^{-1}\sum_s\ind\{B_s\in\{0,1\},A_{s,r}\ne B_s\}$. Both references will be compared with the same released-label interval verdict on the same cohorts; its decisive coverage will also be reported. This measures disagreement, not error against population truth, because $B_s$ itself can be wrong. Six mechanism-specific contrasts will average the two classifier results within each seed. Two-sided exact paired sign tests will assess directional consistency among nonzero seed differences, not equality of mean effects. Ties will be counted, and an all-tie contrast will receive $p=1$. Holm adjustment at family level 0.05 will apply to these six tests. Certified-harm proxy scores and reference-specific native input-state scores will remain descriptive, avoiding a significance claim that silently compares different truth-selected cohorts or different reference targets. Other sensitivities will be exploratory; the identical loss-only control will remain outside hypothesis testing \cite{holm,cerqueira}.

The final budget will be fixed at 30 seeds. Planning will concern interval width rather than a test of the already known nonnegative ordering. Five development seeds will estimate endpoint variability, with substantial uncertainty from that small development sample. The normal-approximation planning width $z_{0.99}s/\sqrt{30}$ will be about 0.021, 0.042 or 0.085 for endpoint standard deviations 0.05, 0.10 or 0.20, respectively. These are design calculations, not experimental precision or power. Oracle identification width will add uncertainty that replication alone cannot remove. The conservative finite-sample envelope will be much wider than this normal approximation; failure to distinguish useful yield from uncertainty will require an inconclusive conclusion. No final-test outcome will trigger extra seeds, a wider harm margin, a different mechanism mixture or a changed primary contrast \cite{hoeffding,souza}.

The statistical report will separate three sources of uncertainty: finite oracle precision, variation across independent seeds and Monte Carlo error from bootstrap resampling. Analysis resampling will use its own recorded random-number seed and preserve the complete paired seed vector. The widths of the identification interval and sampling envelope will be reported separately, together with the exact number of independent seeds. A narrow resampling distribution conditional on one frozen configuration will not establish robustness to different development sets, deployments or calibration choices \cite{hoeffding,souza,nist}.

\subsection*{Step 10 --- Ablation and Component Analysis}
The input-monitor ablation will remove joint MMD and recalibrate the remaining maximum under the same stationary alert target. It may affect input alarms, screened-feature sets and gated positive verdicts, while leaving the evidence-separated loss verdicts identical. The primary matched interval comparison will isolate gating of harm declarations. The strict all-verdict gate will be a separate ablation. Point-loss and point-gated comparisons will examine the effects of removing uncertainty handling and of interpreting absent alarms as below-margin; they will not be presented as additional pure gating contrasts \cite{greco,gretton}.

An uncertainty sensitivity will vary the prespecified interval levels. Window-size, harm-margin and reference sensitivities will preserve observation-level data and seeds. Window-size runs will keep delays, onset times and assessment deadlines fixed in observation units, rather than inadvertently changing the latency treatment by redefining a window. No ablation will infer unavailable labels or introduce adaptive retraining \cite{alencar,cerqueira}.

Component experiments will change one defined element while holding the stream, classifier, risk reference and assessment schedule fixed. Removing MMD will change input evidence but leave the separated loss verdict identical; changing only the harm gate will leave the interval identical; replacing the interval by a point rule will intentionally change uncertainty handling. Sensitivities will vary one factor at a time around the primary configuration unless an interaction is explicitly part of the core crossed design. Confidence, margin, window, amplitude and event-tolerance variants will not automatically form an unaffordable full Cartesian product. A component record will specify the changed element, recalibration requirement, invariant quantities and intended interpretation before that analysis begins \cite{cerqueira,alencar,gretton}.

\subsection*{Step 11 --- Error / Failure Analysis}
The proposed failure catalogue will distinguish harmless input alarms, missed marginal changes, missed dependence changes, false reassurance, unsupported harm declarations, interval overcoverage or undercoverage, prolonged deferral and reference assimilation of changed inputs. Each case will retain the cohort, input-reference membership, label-arrival times, expected-risk truth interval and realized losses. Diagnostics will distinguish deficient evidence from an incorrect decisive statement \cite{greco,szucs,souza,alencar}.

Conditional-only cases will test the limit of input monitoring rather than be used to claim universal input detectability. Dependence-only controls will be checked analytically and by univariate distribution diagnostics. Natural deterioration will not automatically be assigned a certified drift cause. Selective or permanently missing labels, unseen classes, analyst usefulness and optimal retraining action will remain outside the core study's claims \cite{driftvis,souza,alencar}.

Failure analysis will first enumerate every qualifying case and its denominator, then provide the first chronological example of each prespecified failure category within each mechanism, classifier and reference policy. Additional extreme cases may illustrate severity but will be labelled selected examples rather than representative population evidence. Each trace will show current and reference membership, released-label fraction, interval position relative to the margin, oracle bounds when available and the identified failing rule or assumption, with an explicit unknown diagnosis when the available evidence cannot establish one. This will separate a detector's native-task failure, a wrong harm interpretation, legitimate deferral and a data-pipeline fault \cite{souza,window,nist}.

\subsection*{Step 12 --- Final Research Output}
